\documentclass[11pt,a4paper]{article}
\usepackage[utf8]{inputenc}
\usepackage[T1]{fontenc}
\usepackage{lmodern}
\usepackage[french]{babel}
\usepackage{amsmath,amssymb,mathtools}
\usepackage{icomma}
\usepackage{xcolor}
\usepackage{colortbl}
\usepackage{tikz}
\usepackage[most]{tcolorbox}
\usepackage{enumitem}
\usepackage{array}
\usepackage{booktabs}
\usepackage{fancyhdr}
\usepackage[hidelinks]{hyperref}
\usetikzlibrary{arrows.meta,positioning,calc,patterns,backgrounds,shapes.geometric,decorations.pathreplacing,decorations.markings}
\usepackage[a4paper,margin=2.15cm,headheight=26pt,headsep=12pt]{geometry}
\usepackage{titlesec}

% ---------------- Palette ----------------
\definecolor{primary}{RGB}{0,151,169}
\definecolor{primarydark}{RGB}{0,88,101}
\definecolor{defcol}{RGB}{0,114,178}
\definecolor{thmcol}{RGB}{0,140,120}
\definecolor{formulacol}{RGB}{198,93,9}
\definecolor{excol}{RGB}{0,135,90}
\definecolor{attncol}{RGB}{200,40,55}
\definecolor{methcol}{RGB}{90,90,100}
\definecolor{remarkcol}{RGB}{120,94,200}
\definecolor{barcol}{RGB}{0,151,169}
\definecolor{barcold}{RGB}{0,110,124}
\definecolor{modecol}{RGB}{200,55,55}
\definecolor{meancol}{RGB}{0,90,200}
\definecolor{medcol}{RGB}{160,90,190}
\definecolor{quartcol}{RGB}{0,135,90}
\definecolor{ogivecol}{RGB}{176,74,0}

% ---------------- Boîtes ----------------
\tcbset{boxbase/.style={enhanced,breakable,boxrule=0.9pt,arc=2.5pt,
  left=3mm,right=3mm,top=2.3mm,bottom=2.3mm,
  fonttitle=\bfseries,coltitle=white,toptitle=0.6mm,bottomtitle=0.6mm}}

\newtcolorbox{definition}[1][]{boxbase,colback=defcol!6,colframe=defcol,
  title={\textsc{Définition}\ifx\relax#1\relax\relax\else\textmd{ : \itshape #1}\fi}}
\newtcolorbox{propriete}[1][]{boxbase,colback=thmcol!7,colframe=thmcol,
  title={\textsc{Propriété}\ifx\relax#1\relax\relax\else\textmd{ : \itshape #1}\fi}}
\newtcolorbox{formule}[1][]{boxbase,colback=formulacol!7,colframe=formulacol,
  title={\textsc{Formule}\ifx\relax#1\relax\relax\else\textmd{ : \itshape #1}\fi}}
\newtcolorbox{exemple}[1][]{boxbase,colback=excol!6,colframe=excol,
  title={\textsc{Exemple}\ifx\relax#1\relax\relax\else\textmd{ : \itshape #1}\fi}}
\newtcolorbox{methode}[1][]{boxbase,colback=methcol!8,colframe=methcol,sharp corners,
  title={\textsc{Méthode}\ifx\relax#1\relax\relax\else\textmd{ : \itshape #1}\fi}}
\newtcolorbox{remarque}[1][]{boxbase,colback=remarkcol!7,colframe=remarkcol,
  title={\textsc{Astuce}\ifx\relax#1\relax\relax\else\textmd{ : \itshape #1}\fi}}
\newtcolorbox{attention}[1][]{boxbase,colback=attncol!7,colframe=attncol,
  title={\textsc{Attention}\ifx\relax#1\relax\relax\else\textmd{ : \itshape #1}\fi}}

% Boîte "exercice" et "corrigé" (titre libre passé en argument)
\newtcolorbox{exobox}[1]{boxbase,colback=primary!4,colframe=primary,
  before skip=8pt,after skip=8pt,title={#1}}
\newtcolorbox{corbox}[1]{boxbase,colback=excol!5,colframe=excol!85!black,
  before skip=8pt,after skip=8pt,title={#1}}

% ---------------- En-tête ----------------
\newcommand{\docpart}[1]{\def\thedocpart{#1}}
\docpart{}
\pagestyle{fancy}
\fancyhf{}
\fancyhead[L]{\small\textcolor{primarydark}{\textbf{Fiche 8} — Statistiques descriptives}}
\fancyhead[R]{\small\textcolor{primarydark}{\textsc{\thedocpart}}}
\fancyfoot[C]{\small\thepage}
\renewcommand{\headrulewidth}{0.8pt}
\renewcommand{\headrule}{\hbox to\headwidth{\color{primary}\leaders\hrule height \headrulewidth\hfill}}

\titleformat{\section}{\large\bfseries\color{primarydark}}{\thesection}{0.6em}{}
\titleformat{\subsection}{\normalsize\bfseries\color{primary}}{\thesubsection}{0.6em}{}
\setlist{topsep=2pt,itemsep=1.5pt,parsep=0pt}

% Bandeau de titre du document
\newcommand{\bandeau}[2]{%
\begin{tcolorbox}[boxbase,colback=primary,colframe=primarydark,halign=center,
   before skip=2pt,after skip=12pt]
{\LARGE\bfseries\color{white} #1}\\[3pt]{\normalsize\color{white} #2}
\end{tcolorbox}}

\newcommand{\ecm}{\sigma_m}
\newcommand{\ect}{\sigma}
\newcommand{\med}{M\!e}
\docpart{Difficile — Corrigé — Thème 1}
\begin{document}
\bandeau{Thème 1 — Le tableau statistique}{Niveau Difficile — corrigé}

\begin{corbox}{Corrigé 1 — Reconstitution d'un tableau}
\textbf{(a)} $n_1=f_1\times N=0{,}05\times 200=\mathbf{10}$.\\[2pt]
\textbf{(b)} Cumul des deux premières lignes : $n_1+n_2=10+30=40$. Or à la ligne~3 le cumul vaut
$110$, donc $n_3=110-40=\mathbf{70}$.\\[2pt]
\textbf{(c)} $n_4=f_4\times N=0{,}25\times200=\mathbf{50}$.\\[2pt]
\textbf{(d)} $n_5=N-(n_1+n_2+n_3+n_4)=200-(10+30+70+50)=200-160=\mathbf{40}$.\\[2pt]
\textbf{(e)} Tableau complété :
\begin{center}
\renewcommand{\arraystretch}{1.2}
\rowcolors{2}{primary!5}{white}
\begin{tabular}{|c|c|c|c|c|}
\hline
\rowcolor{primary!18}
Classes (min) & $n_i$ & $f_i$ & $n_c$ & $f_c$ \\\hline
$[0,10[$   & 10 & 0{,}05 & 10  & 0{,}05 \\
$[10,20[$  & 30 & 0{,}15 & 40  & 0{,}20 \\
$[20,30[$  & 70 & 0{,}35 & 110 & 0{,}55 \\
$[30,40[$  & 50 & 0{,}25 & 160 & 0{,}80 \\
$[40,50[$  & 40 & 0{,}20 & 200 & 1{,}00 \\\hline
\end{tabular}
\end{center}
Moins de $30$ min : $n_c=\mathbf{110}$ consultations. \; $30$ min ou plus :
$200-110=90$, soit $\dfrac{90}{200}=\mathbf{45\,\%}$.
\end{corbox}

\begin{corbox}{Corrigé 2 — Classes d'amplitudes inégales}
\textbf{(a)} Amplitudes $b-a$ : $10-0=10$ ; $30-10=20$ ; $60-30=30$. Elles \textbf{ne sont pas égales}.\\[2pt]
\textbf{(b)} Centres $x_i=(a+b)/2$ :
\[
x_1=\frac{0+10}{2}=5,\qquad x_2=\frac{10+30}{2}=20,\qquad x_3=\frac{30+60}{2}=45 .
\]
La formule ne fait intervenir que \emph{les deux bornes} : elle donne le milieu quelle que soit
l'amplitude. Elle est donc \textbf{indépendante} de la largeur de la classe.\\[2pt]
\textbf{(c)} $f_1=\dfrac{8}{50}=0{,}16$, $f_2=\dfrac{22}{50}=0{,}44$, $f_3=\dfrac{20}{50}=0{,}40$.
Contrôle : $0{,}16+0{,}44+0{,}40=\mathbf{1{,}00}$. \checkmark
\end{corbox}
\end{document}
